% !TEX root = ../main.tex
\chapter{Atualizações automáticas}
\label{cap:atualizacoes}

O SAPHO evolui rápido, e a AURORA cuida de se manter atual: um atualizador automático embutido verifica, baixa e instala novas versões --- sempre com o seu aval.

\section{Como funciona, do seu ponto de vista}

\begin{enumerate}
  \item \textbf{Verificação silenciosa}: cerca de seis segundos após abrir a IDE, a AURORA consulta o canal oficial de \emph{releases} (o repositório \repo{sapho} no GitHub). Se você já está na última versão, nada acontece --- nenhum diálogo.
  \item \textbf{Oferta}: havendo versão nova, abre a \textbf{janela de atualização}: versão nova e atual, o tamanho do download e o \emph{changelog} bilíngue com o que mudou. Nada é baixado sem você pedir.
  \item \textbf{Download}: ao aceitar, a barra de progresso mostra percentual, MB, velocidade e tempo restante. Você pode continuar trabalhando.
  \item \textbf{Instalação}: com o download pronto, o botão vira \botao{Reiniciar e instalar}. A AURORA fecha, o instalador roda e a IDE reabre na versão nova.
  \item \textbf{Confirmação}: na primeira abertura pós-atualização, um \emph{toast} discreto confirma: \enquote{O SAPHO foi atualizado para a versão X.Y.Z}.
\end{enumerate}

\figurapendente{Janela de atualização}{Com uma versão nova disponível: changelog, tamanho do download e o botão de baixar. Se possível, um segundo print com o download em progresso.}

A integridade do pacote é verificada (\emph{hash} SHA-512 publicado junto com a \emph{release}) antes de instalar; downloads são incrementais quando possível. Em caso de problema --- sem internet, sem espaço em disco, sem permissão --- a janela explica em linguagem clara e a IDE continua funcionando na versão atual.

\begin{info}
\textbf{Dois repositórios, por desenho:} o desenvolvimento da AURORA acontece no repositório \repo{aurora}; as versões estáveis são publicadas no repositório \repo{sapho}, que é o canal que o seu instalador e o atualizador consomem. Você sempre recebe releases estáveis, nunca o estado intermediário do desenvolvimento.
\end{info}

\section{A toolchain acompanha}

Cada instalador carrega, embutida, a versão da toolchain testada com aquela versão da IDE --- os compiladores do YANC, os simuladores, o GTKWave, os moldes de HDL. Atualizar a AURORA atualiza o conjunto todo, em versões que foram validadas juntas. Não há nada para você gerenciar manualmente.

\section{Este manual também se atualiza}

O manual que você lê é parte do ecossistema versionado: cada mudança relevante na AURORA, no YANC ou nos componentes reflete numa edição nova. A versão do manual acompanha a da plataforma.
